\part{A Theory of Verification Residue}

\chapter{Obligations That Survive Verification}
\label{ch:obligations}

A successful run of a proof checker is an event with a precise content: a term has been found to inhabit a type in a stated environment. Almost every claim that is subsequently made on the strength of that event goes beyond this content. It is said that a theorem has been proved, that a published argument has been confirmed, that a question has been answered, or that a result may be relied upon in further work. Each of these claims carries a commitment that the kernel never examined, and the purpose of this chapter is to give those commitments a uniform representation so that they can be counted, tracked, and eventually discharged. The representation is the obligation record, and the collection of obligations that remain unsettled after a verification event is the \emph{verification residue}. The formal definitions are collected in Appendix~\ref{app:formal}. The present chapter motivates them, fixes their intended reading, and shows on small cases how they behave.

\section{What a successful check leaves unsettled}

Consider the polynomial case with which the monograph opens. The stated theorem is that $p(x) = x^3 - x^2 - x + 1$ is nonnegative for $x \ge -1$. The supplied natural-language argument asserts the factorization $p(x) = (x-1)(x+1)^2$, which is incorrect, since $p(x) = (x+1)(x-1)^2$. A language model asked to translate the argument into Lean returned a proof that compiles, having found the correct multiplicities in the course of the translation. The kernel accepts the term. The first baseline obligation, that $\Gamma \vdash \pi : F$, is discharged. Nothing else has been settled by that fact. The statement obligation, that the Lean proposition expresses the stated theorem, is plausible here and can be discharged by inspection, but inspection is a separate act with a separate certificate. The argument obligation, that the Lean proof follows the supplied argument, is false: the supplied argument is invalid, and the formal proof does not formalize it. A faithful translation of that argument would have failed to compile, and it is the success of the translation that conceals the discrepancy.

The example shows that the three baseline obligations of a verification event can take three different statuses at once. The first is discharged, the second can be discharged with additional evidence, and the third is refuted. A report of the form ``the proof was verified'' collapses these into a single value and in doing so asserts something false of the third. The theory developed in this part prevents the collapse by giving each obligation a status of its own and defining the residue as what remains.

\section{The obligation record}

An obligation is a record $o = (\tau, \varphi, \rho, E, D, \sigma)$ whose fields have the following intended readings. The type $\tau$ classifies the obligation. The types that recur in the monograph are \emph{formal} (kernel acceptance of a stated term), \emph{statement} (correspondence of a formal proposition with a specification), \emph{argument} (correspondence of a formal derivation with a natural-language argument), \emph{dependency} (adequacy of a definition, imported lemma, or axiom on which a result relies), \emph{interpretive} (the choice of $\equiv_{\mathcal{I}}$ and of the intended interpretation), \emph{procedural} (properties of the process that produced the artifact, such as the acceptance loop of Chapter~14), and \emph{assumption} (a premise admitted during a transformation and recorded for later discharge). The content $\varphi$ is the proposition the obligation asserts, stated in a content language equipped with a consequence relation. The provenance $\rho$ records where the obligation arose and which act introduced it. The evidence $E$ is the collection of artifacts bearing on the content, including certificates, counterexamples, and expert judgments. The dependencies $D$ record the other obligations on which the establishment of this one relies. The status $\sigma$ is one of five values.

The five statuses are \textsf{open}, \textsf{supported}, \textsf{discharged}, \textsf{refuted}, and \textsf{defeated}. An obligation is \textsf{open} when no evidence has yet been admitted that bears decisively on its content. It is \textsf{supported} when evidence exists that falls short of an accepted certificate, for example a body of test cases or an informal argument that has not been checked. It is \textsf{discharged} when an accepted certificate establishes the content. It is \textsf{refuted} when an accepted certificate establishes the negation of the content. It is \textsf{defeated} when a certificate that once established the content has been overridden, by the discovery of a flaw in the certificate or in something it depends on, without a certificate for the negation having been obtained. The distinction between the last two is substantive. A defeated certificate leaves the content undetermined, whereas a refutation settles it negatively, and conflating them would allow the failure of an argument for a proposition to be reported as the falsity of the proposition.

An obligation whose discharge loses warrant takes the status \textsf{defeated} when a recorded defeater targets its certificate or a dependency, and \textsf{open} when the certificate is retracted without a defeater, so that the earlier evidence is no longer admitted as decisive; the policy is the ledger's, and the historical record keeps both facts. Each obligation carries exactly one status. Questions that might seem to call for a sixth status are handled in other fields. Conflicting evidence, where some admitted evidence favors the content and some opposes it, is recorded as a property of $E$, and the status remains \textsf{open} or \textsf{supported} according to the policy of the ledger until a certificate is accepted. A policy that distinguished a status \textsf{contested} would be a legitimate alternative, and the partition results of Appendix~\ref{app:formal} would then range over six values and not five. What is not permitted is for the status of one obligation to be a set of values, because the residue accounting of the next section depends on the status function being single valued.

\section{Five ways of being unsettled}

The informal vocabulary of verification does not distinguish several quite different reasons why an obligation may fail to be discharged, and the framework is organized so that they can be told apart. A \emph{logical failure} is a case where a derivation is invalid or a statement is false, and the corresponding obligation is refuted. \emph{Semantic uncertainty} is a case where it is not known whether a formal statement means what it was intended to mean; the statement obligation is open, and no amount of kernel checking bears on it. \emph{Missing evidence} is a case where a relevant check has not been carried out, and the obligation is open or supported. \emph{Conflicting evidence} is a case where checks disagree, and it is recorded in $E$. \emph{Defeated certification} is a case where an earlier discharge no longer stands, and the obligation is defeated.

These distinctions are borrowed in structure from the architecture of Adversaria, where claims, evidence, and defeaters are tracked in a public graph, and the application to mathematical formalization is a proposed extension and not a result of the source paper on which the case material of this monograph rests. The borrowing is not decorative: the notion of warranted discharge in Section~\ref{sec:warranted15} below requires exactly the defeater semantics that Adversaria supplies, because a certificate that depends on a defeated assumption must lose its standing without the history of its acceptance being rewritten.

\section{The residue and its partition}

For a finite set $\mathcal{O}$ of obligations attached to an artifact, the residue is the set of obligations whose status is not \textsf{discharged}. Because the status function is single valued, the residue decomposes as the disjoint union
\[
\mathcal{R}(\mathcal{O}) = \mathcal{R}_{\mathsf{open}} \sqcup \mathcal{R}_{\mathsf{supported}} \sqcup \mathcal{R}_{\mathsf{refuted}} \sqcup \mathcal{R}_{\mathsf{defeated}},
\]
which is Proposition~\ref{prop:partition} of the appendix. The statement is definitional and is recorded as such. Its value is that it forces the accounting to be complete: every obligation of an artifact is either discharged or lies in exactly one of four named components. A report that cites only the number of open obligations is incomplete in a specified way, namely that it omits the refuted and defeated components, which are the components that most directly bear on whether the artifact may be relied on.

A residue is not an error. The point is easily lost, because the word suggests a defect. A residue may consist entirely of assumptions that are standard in the field, interpretive choices that are uncontroversial, and statement obligations that every reader would discharge on inspection. What distinguishes a well-run verification from a poorly run one is not the absence of a residue but its legibility: that each member is named, that its status is honestly recorded, and that nothing has been moved out of the residue by relabeling. The residue of a formal proof of a classical theorem in a mature library may be small and benign. The residue of an announced resolution of a major open problem, formalized by an automated translation that has not been audited, may be large, and its size is not a measure of the quality of the underlying mathematics.

\section{Historical and warranted discharge}
\label{sec:warranted15}

Discharge has two readings, and the monograph is explicit about which is in force at any point. A discharge is \emph{historical} when it records that a certificate was accepted by an identified authority under an identified interpretation with an identified set of dependencies. It is \emph{currently warranted} when, in addition, every dependency of the certificate is itself warranted and no defeater is recorded against the certificate. The first is a fact about the past and is never revised. The second is a function of the present ledger and can be lost.

The need for the distinction arises whenever an assumption introduced at an early stage supports certificates at later stages. A formalization may import a library lemma whose statement is later found to differ from its documentation. Every certificate that relied on the lemma remains in the history, but none remains warranted, and the obligations they discharged, unless another certificate for them remains warranted, are no longer warranted-discharged and take one of the other statuses according to the policy of the ledger (Definition~\ref{def:status}). Proposition~\ref{prop:defeasible} of the appendix states this precisely and shows that, for a fixed set of accepted certificates, warrant is antitone in the set of recorded defeaters, while the historical record only accumulates. A verification record that supports only the historical reading gives an inaccurate picture of the evidential state, and one that supports only the warranted reading erases the history that makes revision auditable.

\section{Theorem correctness and proof fidelity}

Two properties are routinely run together, and the obligation framework separates them. \emph{Theorem correctness} is the property that the stated result is true. \emph{Proof fidelity} is the property that a given proof is a faithful formalization of a given argument. A new formal proof of a true theorem may be mathematically valuable without being a faithful formalization of the supplied proof, as in the polynomial example. A faithful formalization of an invalid argument should expose the failure and not silently repair it. The first corresponds to the statement and formal obligations, the second to the argument obligation, and the independence of the two is the subject of the next chapter.

The distinction has a practical consequence for any acceptance procedure that treats compilation as success. If the only test applied to a revised translation is whether it compiles, the procedure counts the repair of an error and the replacement of a correct argument by a different one as the same outcome, which is the mechanism recorded in Section~3.4 of the source and discussed in Chapter~14. The obligation framework makes the omitted test explicit: the argument obligation was not tested, and its status therefore remains what it was before the revision.
